\section*{Bab \ref{ch:g12:curly-arrows} --- Mekanisme Reaksi: Panah Lengkung}

\begin{solution}{exo:g12:curly-arrows:1}
\ce{OH-}: donor (\omterm{def:g10:lewis-and-shape:lone-pair}{pasangan elektron bebas}, muatan negatif); tidak ada
situs akseptor. \ce{H+}: akseptor (sama sekali tidak memiliki \omterm{def:g9:inside-the-atom:nucleus}{elektron}).
\ce{CH2=CH2}: donor (\omterm{def:g10:lewis-and-shape:multiple-bond}{ikatan rangkapnya}). \ce{CH3-Cl}: donor (\omterm{def:g10:lewis-and-shape:lone-pair}{pasangan
elektron bebas} \ce{Cl}), akseptor (karbonnya, $\delta^+$). \ce{H2O}:
donor (\omterm{def:g10:lewis-and-shape:lone-pair}{pasangan elektron bebas} \ce{O}); \omterm{def:g7:atoms-and-molecules:atom}{atom-atom} hidrogennya,
$\delta^+$, adalah situs akseptor yang lemah.
\end{solution}

\begin{solution}{exo:g12:curly-arrows:2}
(a) \omterm{def:g12:curly-arrows:categories}{substitusi}; (b) \omterm{def:g12:curly-arrows:categories}{adisi}; (c) \omterm{def:g12:curly-arrows:categories}{eliminasi}; (d) \omterm{def:g12:curly-arrows:categories}{adisi}.
\end{solution}

\begin{solution}{exo:g12:curly-arrows:3}
Panah itu berawal dari sepasang \omterm{def:g9:inside-the-atom:nucleus}{elektron} (\omterm{def:g10:lewis-and-shape:lone-pair}{pasangan elektron bebas} atau
sebuah ikatan) dan berakhir pada \omterm{def:g7:atoms-and-molecules:atom}{atom} tempat pasangan itu pergi. Dari
sepasang \omterm{def:g9:inside-the-atom:nucleus}{elektron} bebas oksigen ke sebuah karbon: pasangan itu menjadi
ikatan \ce{O-C} yang baru.
\end{solution}

\begin{solution}{exo:g12:curly-arrows:4}
Amonia memberikan \omterm{def:g10:lewis-and-shape:lone-pair}{pasangan elektron bebas} \omterm{def:g7:atoms-and-molecules:atom}{atom} nitrogennya; panahnya
berawal dari pasangan itu ke \ce{H+}. Produk \ce{NH4+} bermuatan $+1$.
\end{solution}

\begin{solution}{exo:g12:curly-arrows:5}
Spesies yang terbentuk dalam satu tahap dan habis terpakai dalam tahap
berikutnya; spesies itu tidak ada dalam persamaan keseluruhan. Di sini,
karbokation \ce{CH3-CH2+}.
\end{solution}

\begin{solution}{exo:g12:curly-arrows:6}
Satu panah dari sepasang \omterm{def:g9:inside-the-atom:nucleus}{elektron} bebas oksigen \ce{OH-} ke karbon; satu
panah dari ikatan \ce{C-Br} ke brom. Muatan: $-1$ di kiri (\ce{OH-}); di
kanan, metanol netral dan \ce{Br-} bermuatan $-1$: totalnya $-1$ di
kedua ruas.
\end{solution}

\begin{solution}{exo:g12:curly-arrows:7}
Brom lebih elektronegatif daripada hidrogen: brom mengambil pasangan itu
dan lepas sebagai \ce{Br-} ($-1$); hidrogen berikatan dengan sebuah
karbon, dan karbon yang lain tinggal bermuatan positif: \ce{CH3-CH2+}
($+1$).
\end{solution}

\begin{solution}{exo:g12:curly-arrows:8}
1.~\ce{CH3-CH2-OH + H+ -> CH3-CH2-OH2+}: donornya sepasang \omterm{def:g9:inside-the-atom:nucleus}{elektron}
bebas oksigen; akseptornya \ce{H+}. 2.~\ce{CH3-CH2-OH2+ -> CH3-CH2+ + H2O}:
pasangan \ce{C-O} pergi ke oksigen yang positif (akseptor), yang lepas
bersamanya sebagai air. 3.~\ce{CH3-CH2+ -> CH2=CH2 + H+}: \omterm{def:g10:lewis-and-shape:lone-pair}{pasangan
elektron ikatan} \ce{C-H} (donor) berpindah ke karbon yang positif
(akseptor) untuk membentuk \omterm{def:g10:lewis-and-shape:multiple-bond}{ikatan rangkap}, dan \ce{H+} lepas. \omterm{def:g9:ions:ion}{Ion}
hidrogen yang dipakai pada tahap 1 dikembalikan pada tahap 3: \omterm{def:g9:ions:ion}{ion} itu
muncul di kedua ruas dan saling meniadakan.
\end{solution}

\begin{solution}{exo:g12:curly-arrows:9}
Panah dari sepasang \omterm{def:g9:inside-the-atom:nucleus}{elektron} bebas oksigen air ke karbon yang positif.
Setelah \ce{H+} lepas, terbentuk etanol \ce{CH3-CH2-OH}. Keseluruhannya:
\ce{CH2=CH2 + H2O -> CH3-CH2-OH}, sebuah \omterm{def:g12:curly-arrows:categories}{adisi}.
\end{solution}

\begin{solution}{exo:g12:curly-arrows:10}
Dua pasang \omterm{def:g9:inside-the-atom:nucleus}{elektron} berpindah, dalam satu tahap: ikatan \ce{C-O}
terbentuk, ikatan \ce{C-Br} putus. Tidak ada zat antara.
\end{solution}

\begin{solution}{exo:g12:curly-arrows:11}
Klor lebih elektronegatif ($3.16 - 2.55 = 0.61$ dibandingkan
$2.96 - 2.55 = 0.41$): karbon klorometana lebih $\delta^+$, yang dengan
sendirinya akan membuatnya bereaksi lebih cepat. Kenyataannya tidak:
kecepatannya juga bergantung pada seberapa mudah \omterm{def:g10:electron-shells:family}{halogennya} lepas
bersama pasangan \omterm{def:g9:inside-the-atom:nucleus}{elektronnya}, dan brom, yang lebih besar, lebih mudah
lepas.
\end{solution}

\begin{solution}{exo:g12:curly-arrows:12}
Jika hidrogen berikatan dengan karbon 1: \ce{CH3-CH+-CH3}, yang
menghasilkan 2-kloropropana \ce{CH3-CHCl-CH3}. Jika berikatan dengan
karbon 2: \ce{+CH2-CH2-CH3}, yang menghasilkan 1-kloropropana
\ce{CH2Cl-CH2-CH3}.
\end{solution}

\begin{solution}{exo:g12:curly-arrows:13}
Arah panahnya terbalik: panah harus berawal dari donor, yaitu sepasang
\omterm{def:g9:inside-the-atom:nucleus}{elektron} bebas \ce{OH-}, dan berakhir pada karbon akseptor. Selain itu,
ikatan baru terbentuk antara oksigen dan karbon, bukan antara oksigen
dan brom: brom lepas, dengan membawa \omterm{def:g10:lewis-and-shape:lone-pair}{pasangan elektron ikatan}
\ce{C-Br}.
\end{solution}

\begin{solution}{exo:g12:curly-arrows:14}
Karbon gugus \ce{C=O} sebuah asam terikat pada dua \omterm{def:g7:atoms-and-molecules:atom}{atom} oksigen yang
elektronegatif: sangat $\delta^+$, sebuah akseptor. Oksigen \omterm{def:g11:functional-groups:alcohol}{alkohol}
bermuatan $\delta^-$ dan memiliki dua \omterm{def:g10:lewis-and-shape:lone-pair}{pasangan elektron bebas}: sebuah
donor. Secara keseluruhan air dilepaskan; gugus \ce{-OH} asam digantikan
oleh gugus \ce{-O-R} \omterm{def:g11:functional-groups:alcohol}{alkohol}: sebuah \omterm{def:g12:curly-arrows:categories}{substitusi}.
\end{solution}

\begin{solution}{exo:g12:curly-arrows:15}
Karena didorong oleh \omterm{def:g9:inside-the-atom:nucleus}{elektron-elektron} \omterm{def:g10:lewis-and-shape:multiple-bond}{ikatan rangkap}, \omterm{def:g10:lewis-and-shape:lone-pair}{pasangan elektron
ikatan} \ce{Br-Br} berpindah ke arah brom yang jauh: brom yang dekat
menjadi $\delta^+$, sebuah situs akseptor. Tahap pertama: panah dari
\omterm{def:g10:lewis-and-shape:multiple-bond}{ikatan rangkap} ke brom yang dekat, dan panah dari ikatan \ce{Br-Br} ke
brom yang jauh, yang lepas sebagai \ce{Br-}; muatan positif tertinggal
pada karbon lain dari bekas \omterm{def:g10:lewis-and-shape:multiple-bond}{ikatan rangkap} itu.
\end{solution}

\begin{solution}{pb:g12:curly-arrows:1}
\textbf{1.} Etena: \ce{H2C=CH2}, satu \omterm{def:g10:lewis-and-shape:multiple-bond}{ikatan rangkap} dan empat
\ce{C-H}. Hidrogen bromida: \ce{H-Br} dengan tiga \omterm{def:g10:lewis-and-shape:lone-pair}{pasangan elektron
bebas} pada brom.

\textbf{2.} \omterm{def:g10:lewis-and-shape:multiple-bond}{Ikatan rangkapnya}: pasangan \omterm{def:g9:inside-the-atom:nucleus}{elektron} keduanya merupakan
daerah yang kaya \omterm{def:g9:inside-the-atom:nucleus}{elektron}.

\textbf{3.} $2.96 - 2.20 = 0.76$: \ce{H} bermuatan $\delta^+$.

\textbf{4.} \omterm{def:g7:atoms-and-molecules:atom}{Atom} hidrogennya.

\textbf{5.} Panah dari \omterm{def:g10:lewis-and-shape:multiple-bond}{ikatan rangkap} ke hidrogen \ce{HBr}; panah dari
ikatan \ce{H-Br} ke brom.

\textbf{6.} Karbokation \ce{CH3-CH2+} ($+1$) dan \ce{Br-} ($-1$).

\textbf{7.} Sebelum: $0 + 0 = 0$; sesudah: $+1 - 1 = 0$.

\textbf{8.} Panah dari sepasang \omterm{def:g9:inside-the-atom:nucleus}{elektron} bebas \ce{Br-} ke karbon yang
positif.

\textbf{9.} Enam (tiga ikatan): kurang satu pasang dari oktet, sehingga
karbon itu menerima sepasang \omterm{def:g9:inside-the-atom:nucleus}{elektron} dari donor mana pun di dekatnya.

\textbf{10.} Karbokation \ce{CH3-CH2+}.

\textbf{11.} Zat antara itu terbentuk pada tahap pertama dan habis
terpakai pada tahap kedua.

\textbf{12.} Sebuah \omterm{def:g12:curly-arrows:categories}{adisi}.

\textbf{13.} Hanya gugusnya: rantai dua karbon dipertahankan, dan \omterm{def:g10:lewis-and-shape:multiple-bond}{ikatan
rangkapnya} menjadi ikatan tunggal yang membawa \ce{H} dan \ce{Br}.

\textbf{14.} \qty{28.0}{g/mol}; \qty{108.9}{g/mol}.

\textbf{15.} $10.0 / 28.0 = \qty{0.357}{mol}$.

\textbf{16.} Etena: $0.357 - x$; \ce{HBr}: berlebih; bromoetana: $x$.
Etena menjadi pembatas: $x_{\max} = \qty{0.357}{mol}$.

\textbf{17.} $0.357 \times 108.9 = \qty{38.9}{g}$.

\textbf{18.} $0.75 \times 38.9 = \qty{29.2}{g}$.

\textbf{19.} Sekitar \qty{29}{g} bromoetana.
\end{solution}
